\VerdierPhysicalPageAnchor{0071}%
\VerdierSemanticAnchor{ch1-2-3}
\subsection{Structure graduée sur les catégories de complexes}

\VerdierSemanticAnchor{ch1-2-3-1}
\noindent\textbf{2.3.1.}\quad
Soient $Y$ un complexe $n$-uple et $i\in[n]$. Posons :

\VerdierSemanticAnchor{ch1-2-3-1-eq1}
\begin{equation}
  \begin{aligned}
    (T_iY)^\sigma&=Y^{\sigma+e_i},
      &&\sigma\in\mathbb Z[n],\\
    d_{T_iY}^{\sigma,j}&=\nu(i,j)d_Y^{\sigma+e_i,j},
      &&j\in[n],
  \end{aligned}
  \tag{2.3.1.1}
\end{equation}
où $\nu(i,j)=1$ si $i\ne j$, $\nu(i,j)=-1$ si $i=j$.

\VerdierSemanticAnchor{ch1-2-3-1-morphism-shift}
Pour tout morphisme de complexes $n$-uples
$f\colon Y\longrightarrow Y'$, posons :

\VerdierSemanticAnchor{ch1-2-3-1-eq2}
\begin{equation}
  (T_if)^\sigma=f^{\sigma+e_i},
  \qquad \sigma\in\mathbb Z[n].
  \tag{2.3.1.2}
\end{equation}

\VerdierSemanticAnchor{ch1-2-3-1-translation-functor}
On définit ainsi un foncteur
$T_i\colon\operatorname{comp}^{n}(C)\longrightarrow
\operatorname{comp}^{n}(C)$ qu'on appelle le
\emph{foncteur de translation dans la direction $i$}.

\VerdierSemanticAnchor{ch1-2-3-1-n-equals-one}
Lorsque $n=1$, on utilise la notation $TY$, $Tf$ ou bien $Y[1]$, $f[1]$.
Le foncteur de translation dans la direction $i$ est un automorphisme de
$\operatorname{comp}^{n}(C)$.

\VerdierSemanticAnchor{ch1-2-3-1-eq3}
Les foncteurs $T_i$, $i\in[n]$, commutent deux à deux :
\begin{equation}
  T_iT_j=T_jT_i,
  \qquad \forall i,j\in[n].
  \tag{2.3.1.3}
\end{equation}

\VerdierSemanticAnchor{ch1-2-3-1-action}
On peut donc définir une opération du groupe $\mathbb Z[n]$ sur la catégorie
$\operatorname{comp}^{n}(C)$ :

\VerdierSemanticAnchor{ch1-2-3-1-eq4}
\begin{equation}
  \sigma=\sum_{i\in[n]}l_ie_i,
  \qquad
  T(\sigma)=T_1^{l_1}\cdots T_n^{l_n}.
  \tag{2.3.1.4}
\end{equation}

\VerdierSemanticAnchor{ch1-2-3-1-strict-graded-category}
La notation $\operatorname{comp}^{n}(C)$ désignera désormais la catégorie
$\operatorname{comp}^{n}(C)$ munie de la structure de
$\mathbb Z[n]$-catégorie stricte décrite ci-dessus.

\VerdierSemanticAnchor{ch1-2-3-2}
\noindent\textbf{2.3.2.}\quad
Soient $\varphi\colon[n]\longrightarrow[m]$ une application, $Y$ un complexe
$n$-uple de $C$, $i$ un élément de $[n]$. On suppose, pour fixer les idées,
que les sommes dénombrables sont représentables dans $C$.

\VerdierSemanticAnchor{ch1-2-3-2-equal-objects}
Remarquons tout d'abord que les objets
$\left(\int_\varphi^\Sigma T_iY\right)^\tau$ et
$\left(T_{\varphi(i)}\int_\varphi^\Sigma Y\right)^\tau$ sont égaux.
Ils sont en effet égaux par définition à :

\VerdierSemanticAnchor{ch1-2-3-2-direct-sum}
\begin{equation*}
  \bigoplus_{\varphi(\sigma)=\tau+e_{\varphi(i)}}Y^\sigma.
\end{equation*}
Posons alors :

\VerdierSemanticAnchor{ch1-2-3-2-eq1}
\begin{equation}
  \begin{aligned}
    \bigl(p_{i,\varphi}(Y)\bigr)^\tau\colon
    \left(\int_\varphi^\Sigma T_iY\right)^\tau
      &\longrightarrow
    \left(T_{\varphi(i)}\int_\varphi^\Sigma Y\right)^\tau,
      &&\tau\in\mathbb Z[m],\\[2pt]
    \bigl(p_{i,\varphi}(Y)\bigr)^\tau
      &=\bigoplus_{\varphi(\sigma)=\tau+e_{\varphi(i)}}
        \varepsilon(\sigma,i)\varepsilon(\varphi(\sigma),\varphi(i))
        \operatorname{id}_{Y^\sigma}.
  \end{aligned}
  \tag{2.3.2.1}
\end{equation}
